\documentclass[ ../main.tex]{subfiles}
\providecommand{\mainx}{..}
\begin{document}
\section{Introduction}
\label{sec:intro}

Bernoulli sets~\cite{bernoulliSets} formalize approximate set membership: the indicator function $\SetIndicator{\Set{A}} \colon \Set{U} \to \{0,1\}$ is approximated with independent per-element errors.
Bernoulli maps~\cite{bernoulliMaps} generalize this to arbitrary functions $\Fun{f} \colon \Set{X} \to \Set{Y}$, where the per-input errors are independent and the joint confusion matrix factors as a Kronecker product of per-input channels.
But real programs compose \emph{types}, not just sets and functions.
A database record is a product type (tuples of fields).
A nullable value is a sum type (the value or nothing).
A query result is a function type (input to output).
When these types carry approximation error, how does that error propagate through the type constructors?

This paper provides the constructive counterpart to the function-space viewpoint developed in the companion paper on Bernoulli maps~\cite{bernoulliMaps}.
Where that paper takes a \emph{denotational} perspective (every approximation is a latent function approximated per-input), this paper takes a \emph{constructive} perspective: we build approximate types from primitive types using algebraic type constructors (sum, product, exponential) and derive how approximation distributes over each constructor.

\paragraph{Central claim.}
Every algebraic data type constructor has a well-defined random approximate counterpart with quantifiable error propagation.
Product types are well-behaved: the joint confusion matrix factors as a Kronecker product, and parameters grow additively.
Sum types introduce a structural complication: tag errors couple the components and break the factorization.
Exponential types (function spaces) are precisely Bernoulli maps.
The two perspectives (denotational and constructive) converge at the exponential type, yielding the same Kronecker factorization theorem.

\paragraph{Contributions.}
\begin{enumerate}
\item We define the random approximate value $\AT{T}$ for each algebraic type constructor: void, unit, Bool, sum, product, and exponential (\cref{sec:prims,sec:composite}).

\item We prove that product types preserve the Kronecker factorization (\cref{thm:product_parsimony}): an approximate product $\AT{(\Set{A} \times \Set{B})}$ has $\mathrm{params}(A) + \mathrm{params}(B)$ free parameters, not the exponentially larger count of a general channel.

\item We show that sum types break the factorization (\cref{rem:sum_nonfactor}): $\AT{(\Set{A} + \Set{B})} \not\cong \AT{\Set{A}} + \AT{\Set{B}}$ because tag errors change which component's error model applies.

\item We identify the optional type $\mathrm{Maybe}(\Set{A}) = \Set{A} + \Unit$ as the type-theoretic formulation of Bernoulli set membership, with tag errors corresponding to false positives and false negatives.
\end{enumerate}

\paragraph{Organization.}
\Cref{sec:approx_values} defines the random approximate value and develops the approximate bit as the canonical worked example.
\Cref{sec:prims} defines the primitive approximate types (void, unit, Bool).
\Cref{sec:composite} derives the composite type constructors (product, sum, exponential) and proves the Kronecker factorization theorem for products.
\Cref{sec:error_prop} shows how errors propagate through function application, with a worked example for the approximate AND gate.
\Cref{sec:bernoulli_boolean_algebra} develops the approximate Boolean algebra over bit vectors.
\Cref{sec:invariants} addresses data types with invariants and quantifies the probability of invariant violation.
\Cref{sec:cpp} provides a computational realization as a C++ monad.

\end{document}
